\NativeSlot{48}{65}{%
% p071_s01

\ExamQuestion{7.}{5}%
下列交通标志，可以经过旋转得到的是\mbox{（\hspace{1.8em}）}。选A、B、C或D。\par
\BookFigure{figures/p071_s01-01.tex}
}
\NativeSlot{160}{156}{%
\ExamGroup{四、画一画。{\normalfont\kaishu\fontsize{10.8bp}{14bp}\selectfont （共18分）}}

% p071_s02

\ExamQuestion{1.}{5}%
下图中，点O是线段AB上的一点，将线段AB绕点O顺时针方向旋转\(90^\circ\)，旋转后的线段与原线段的位置有何关系？如果逆时针方向旋转\(90^\circ\)呢？请分别在下面的图上画出来。\par
\BookFigure{figures/p071_s02-01.tex}

% p071_s03

\ExamQuestion{2.}{5}%
按要求画图。\par\noindent （1）画出图①的另一半，使它成为一个轴对称图形。\par

% p071_s04

\ExamQuestion{（2）}{5}%
画出图②向下平移7格后的图形。\par

% p071_s05

\ExamQuestion{（3）}{5}%
画出图③绕点A逆时针旋转\(90^\circ\)后的图形。\par
\BookFigure{figures/p071_s03-01.tex}
}
\NativeSlot{615}{106}{%
\par\noindent%
\ExamTile{.485}{%
% p071_s06

\ExamQuestion{3.}{5}%
围绕点O，旋转三角形画一个风车。\par
\BookFigure{figures/p071_s06-01.tex}
}\hfill%
\ExamTile{.485}{%
% p071_s07

\ExamQuestion{4.}{5}%
画出镜子中看到的时刻。原钟的分针指向12，时针指向4。\par
\BookFigure{figures/p071_s07-01.tex}
}\par
}
